% GENERATED FILE. Edit canonical lessons, then run npm run generate:books.
% canonical-source-hash: fnv1a64:3eaef5cea971c086
% canonical-lessons: BN-C299-somprochar, BN-C299-chenel, BN-C299-uposthapok, BN-C299-sorasori, BN-C299-dorshok, BN-R299-review-292-299

\chapter{On the Air}
\label{ch:bn-on-the-air}
\hlchaptermodality{299}

\begin{chapteropening}
\textbf{By the end of this chapter:} \emph{I can talk about broadcasts, channels and live news.}

Complete the last lesson of chapter 299: I can talk about broadcasts, channels and live news.
\end{chapteropening}

\section[\texorpdfstring{\bn{সম্প্রচার} (sômprôchār)}{সম্প্রচার (sômprôchār)}]{\bn{সম্প্রচার} (sômprôchār) — a broadcast}

\label{lesson:BN-C299-somprochar}



\begin{quote}
Before the new one: say the Bengali for to whisper, then the Bengali for a sensation.
\end{quote}

\subsection*{You'll want to know: \bn{সম্প্রচার}}
\textbf{\bn{সম্প্রচার}} — \emph{sômprôchār} — \textquotedblleft{}a broadcast\textquotedblright{}.

\begin{quote}
\bn{রেডিওতে} \bn{সকাল} \bn{আটটায়} \bn{খবরের} \bn{সম্প্রচার} \bn{হয়।}
\end{quote}

\begin{quote}
\emph{reḍiote sôkāl āṭṭāy khôbôrer sômprôchār hôy.}
\end{quote}

\begin{quote}
The news is broadcast on the radio at eight in the morning.
\end{quote}

\begin{cousinweb}
\textbf{\bn{সম্প্রচার}} (\emph{sômprôchār}) is a learned word built from Sanskrit \emph{sam-}, 'together', and \emph{pracāra}, 'a going forth, a spreading' (Monier-Williams): a message sent out to everyone at once.
\end{cousinweb}

\begin{grammarlens}[title={What you've built}]
One more word for telling what happened.
\end{grammarlens}

\subsection*{Your turn}
\begin{itemize}
\raggedright
  \item \emph{Say it:} \emph{sômprôchār}
  \item \emph{Say it:} \emph{sômprôchār}, once more
  \item \emph{Say it:} \emph{reḍiote sôkāl āṭṭāy khôbôrer sômprôchār hôy.} — The news is broadcast on the radio at eight in the morning.
  \item \emph{Recall:} say the Bengali for to whisper, then the Bengali for a sensation, then say \emph{sômprôchār} again
\end{itemize}

\begin{tcolorbox}[breakable,colback=teal!4,colframe=teal!35!black,title={Before you move on}]
What does \bn{সম্প্রচার} mean? (\textquotedblleft{}A broadcast\textquotedblright{}.) Say it once more.
\end{tcolorbox}
\section[\texorpdfstring{\bn{চ্যানেল} (chænel)}{চ্যানেল (chænel)}]{\bn{চ্যানেল} (chænel) — a channel}

\label{lesson:BN-C299-chenel}



\begin{quote}
Before the new one: say the Bengali for a sensation, then the Bengali for a broadcast.
\end{quote}

\subsection*{You'll want to know: \bn{চ্যানেল}}
\textbf{\bn{চ্যানেল}} — \emph{chænel} — \textquotedblleft{}a channel\textquotedblright{}.

\begin{quote}
\bn{এই} \bn{চ্যানেলে} \bn{সারাদিন} \bn{খবর} \bn{দেখায়।}
\end{quote}

\begin{quote}
\emph{ei chænele sārādin khôbôr dekhāy.}
\end{quote}

\begin{quote}
This channel shows news all day.
\end{quote}

\begin{cousinweb}
\textbf{\bn{চ্যানেল}} (\emph{chænel}) is borrowed from English \emph{channel}. The \bn{চ্যা} at the start is read \emph{æ}, as in \textbf{\bn{ব্যাটারি}}.
\end{cousinweb}

\begin{grammarlens}[title={What you've built}]
One more word for telling what happened.
\end{grammarlens}

\subsection*{Your turn}
\begin{itemize}
\raggedright
  \item \emph{Say it:} \emph{chænel}
  \item \emph{Say it:} \emph{chænel}, once more
  \item \emph{Say it:} \emph{ei chænele sārādin khôbôr dekhāy.} — This channel shows news all day.
  \item \emph{Recall:} say the Bengali for a sensation, then the Bengali for a broadcast, then say \emph{chænel} again
\end{itemize}

\begin{tcolorbox}[breakable,colback=teal!4,colframe=teal!35!black,title={Before you move on}]
What does \bn{চ্যানেল} mean? (\textquotedblleft{}A channel\textquotedblright{}.) Say it once more.
\end{tcolorbox}
\section[\texorpdfstring{\bn{উপস্থাপক} (upôsthāpôk)}{উপস্থাপক (upôsthāpôk)}]{\bn{উপস্থাপক} (upôsthāpôk) — a presenter; a host}

\label{lesson:BN-C299-uposthapok}



\begin{quote}
Before the new one: say the Bengali for a broadcast, then the Bengali for a channel.
\end{quote}

\subsection*{You'll want to know: \bn{উপস্থাপক}}
\textbf{\bn{উপস্থাপক}} — \emph{upôsthāpôk} — \textquotedblleft{}a presenter; a host\textquotedblright{}.

\begin{quote}
\bn{উপস্থাপক} \bn{ঝড়ের} \bn{খবর} \bn{পড়লেন।}
\end{quote}

\begin{quote}
\emph{upôsthāpôk jhôṛer khôbôr pôṛlen.}
\end{quote}

\begin{quote}
The presenter read the news about the storm.
\end{quote}

\begin{cousinweb}
\textbf{\bn{উপস্থাপক}} (\emph{upôsthāpôk}) is Sanskrit \emph{upasthāpaka}, 'causing to turn one's attention, reminding' (Monier-Williams); the sense 'a presenter', the one who puts the news before you, is the modern one.
\end{cousinweb}

\begin{grammarlens}[title={What you've built}]
One more word for telling what happened.
\end{grammarlens}

\subsection*{Your turn}
\begin{itemize}
\raggedright
  \item \emph{Say it:} \emph{upôsthāpôk}
  \item \emph{Say it:} \emph{upôsthāpôk}, once more
  \item \emph{Say it:} \emph{upôsthāpôk jhôṛer khôbôr pôṛlen.} — The presenter read the news about the storm.
  \item \emph{Recall:} say the Bengali for a broadcast, then the Bengali for a channel, then say \emph{upôsthāpôk} again
\end{itemize}

\begin{tcolorbox}[breakable,colback=teal!4,colframe=teal!35!black,title={Before you move on}]
What does \bn{উপস্থাপক} mean? (\textquotedblleft{}A presenter; a host\textquotedblright{}.) Say it once more.
\end{tcolorbox}
\section[\texorpdfstring{\bn{সরাসরি} (sôrāsôri)}{সরাসরি (sôrāsôri)}]{\bn{সরাসরি} (sôrāsôri) — live; directly}

\label{lesson:BN-C299-sorasori}



\begin{quote}
Before the new one: say the Bengali for a channel, then the Bengali for a presenter.
\end{quote}

\subsection*{You'll want to know: \bn{সরাসরি}}
\textbf{\bn{সরাসরি}} — \emph{sôrāsôri} — \textquotedblleft{}live; directly\textquotedblright{}.

\begin{quote}
\bn{টিভিতে} \bn{খেলাটা} \bn{সরাসরি} \bn{দেখাল।}
\end{quote}

\begin{quote}
\emph{ṭibhite khelāṭā sôrāsôri dekhālô.}
\end{quote}

\begin{quote}
The match was shown live on television.
\end{quote}

\begin{cousinweb}
\textbf{\bn{সরাসরি}} (\emph{sôrāsôri}) is usually traced to Persian \emph{sarāsar}, 'from end to end, throughout'; Bengali uses it for 'straight, directly' and for a broadcast shown live.
\end{cousinweb}

\begin{grammarlens}[title={What you've built}]
One more word for telling what happened.
\end{grammarlens}

\subsection*{Your turn}
\begin{itemize}
\raggedright
  \item \emph{Say it:} \emph{sôrāsôri}
  \item \emph{Say it:} \emph{sôrāsôri}, once more
  \item \emph{Say it:} \emph{ṭibhite khelāṭā sôrāsôri dekhālô.} — The match was shown live on television.
  \item \emph{Recall:} say the Bengali for a channel, then the Bengali for a presenter, then say \emph{sôrāsôri} again
\end{itemize}

\begin{tcolorbox}[breakable,colback=teal!4,colframe=teal!35!black,title={Before you move on}]
What does \bn{সরাসরি} mean? (\textquotedblleft{}Live; directly\textquotedblright{}.) Say it once more.
\end{tcolorbox}
\section[\texorpdfstring{\bn{দর্শক} (dôrshôk)}{দর্শক (dôrshôk)}]{\bn{দর্শক} (dôrshôk) — a viewer; a spectator}

\label{lesson:BN-C299-dorshok}



\begin{quote}
Before the new one: say the Bengali for a presenter, then the Bengali for live.
\end{quote}

\subsection*{You'll want to know: \bn{দর্শক}}
\textbf{\bn{দর্শক}} — \emph{dôrshôk} — \textquotedblleft{}a viewer; a spectator\textquotedblright{}.

\begin{quote}
\bn{খেলাটা} \bn{দেখতে} \bn{স্টেডিয়ামে} \bn{হাজার} \bn{হাজার} \bn{দর্শক} \bn{এল।}
\end{quote}

\begin{quote}
\emph{khelāṭā dekhte sṭeḍiyāme hājār hājār dôrshôk elô.}
\end{quote}

\begin{quote}
Thousands of spectators came to the stadium to watch the match.
\end{quote}

\begin{cousinweb}
\textbf{\bn{দর্শক}} (\emph{dôrshôk}) is Sanskrit \emph{darśaka}, 'seeing, one who looks' (Monier-Williams), from \emph{√dṛś}, 'to see'.
\end{cousinweb}

\begin{grammarlens}[title={What you've built}]
One more word for telling what happened.
\end{grammarlens}

\subsection*{Your turn}
\begin{itemize}
\raggedright
  \item \emph{Say it:} \emph{dôrshôk}
  \item \emph{Say it:} \emph{dôrshôk}, once more
  \item \emph{Say it:} \emph{khelāṭā dekhte sṭeḍiyāme hājār hājār dôrshôk elô.} — Thousands of spectators came to the stadium to watch the match.
  \item \emph{Recall:} say the Bengali for a presenter, then the Bengali for live, then say \emph{dôrshôk} again
\end{itemize}

\begin{tcolorbox}[breakable,colback=teal!4,colframe=teal!35!black,title={Before you move on}]
What does \bn{দর্শক} mean? (\textquotedblleft{}A viewer; a spectator\textquotedblright{}.) Say it once more.
\end{tcolorbox}
\section[(dialogue)]{Looking back: from Drought and Hunger to On the Air}

\label{lesson:BN-R299-review-292-299}



\begin{quote}
Say the Bengali for a drought and a famine.
\end{quote}

\subsection*{Your turn}
Say these aloud:

\begin{itemize}
\raggedright
  \item a drought, a famine, starvation, to crack, to be scorched
  \item an earthquake, a tremor, to collapse, destruction, to be trapped underneath
  \item alert, careful, a warning sign, an apprehension, to avoid
  \item damage, rubble, to mend, to tackle, to regain
  \item a newspaper, a headline, a report, to publish, an editor
  \item talk, an interview, to announce, a public statement, to mention
  \item tension, to spread a rumour, an uproar, to whisper, a sensation
  \item a broadcast, a channel, a presenter, live, a viewer
\end{itemize}

\begin{tcolorbox}[breakable,colback=teal!4,colframe=teal!35!black,title={Before you move on}]
A drought? (\textbf{\bn{খরা}}, \emph{khôrā}.) A famine? (\textbf{\bn{দুর্ভিক্ষ}}, \emph{durbhikkhô}.) Starvation? (\textbf{\bn{অনাহার}}, \emph{ônāhār}.) To crack? (\textbf{\bn{ফাটা}}, \emph{phāṭā}.) To be scorched? (\textbf{\bn{ঝলসে} \bn{যাওয়া}}, \emph{jhôlse jāoyā}.) An earthquake? (\textbf{\bn{ভূমিকম্প}}, \emph{bhūmikômpô}.) A tremor? (\textbf{\bn{কম্পন}}, \emph{kômpôn}.) To collapse? (\textbf{\bn{ধসে} \bn{পড়া}}, \emph{dhôse pôṛā}.) Destruction? (\textbf{\bn{ধ্বংস}}, \emph{dhvôngsô}.) To be trapped underneath? (\textbf{\bn{চাপা} \bn{পড়া}}, \emph{chāpā pôṛā}.) Alert? (\textbf{\bn{সতর্ক}}, \emph{sôtôrkô}.) Careful? (\textbf{\bn{সাবধান}}, \emph{sābdhān}.) A warning sign? (\textbf{\bn{সংকেত}}, \emph{sôngket}.) An apprehension? (\textbf{\bn{আশঙ্কা}}, \emph{āshôngkā}.) To avoid? (\textbf{\bn{এড়ানো}}, \emph{eṛāno}.) Damage? (\textbf{\bn{ক্ষতি}}, \emph{khôti}.) Rubble? (\textbf{\bn{ধ্বংসস্তূপ}}, \emph{dhvôngsôstūp}.) To mend? (\textbf{\bn{মেরামত} \bn{করা}}, \emph{merāmôt kôrā}.) To tackle? (\textbf{\bn{মোকাবিলা} \bn{করা}}, \emph{mokābilā kôrā}.) To regain? (\textbf{\bn{ফিরে} \bn{পাওয়া}}, \emph{phire pāoyā}.) A newspaper? (\textbf{\bn{সংবাদপত্র}}, \emph{sôngbādpôtrô}.) A headline? (\textbf{\bn{শিরোনাম}}, \emph{shironām}.) A report? (\textbf{\bn{প্রতিবেদন}}, \emph{prôtibedôn}.) To publish? (\textbf{\bn{প্রকাশ} \bn{করা}}, \emph{prôkāsh kôrā}.) An editor? (\textbf{\bn{সম্পাদক}}, \emph{sômpādôk}.) Talk? (\textbf{\bn{কথা}}, \emph{kôthā}.) An interview? (\textbf{\bn{সাক্ষাৎকার}}, \emph{sākkhātkār}.) To announce? (\textbf{\bn{ঘোষণা} \bn{করা}}, \emph{ghoshônā kôrā}.) A public statement? (\textbf{\bn{বিবৃতি}}, \emph{bibriti}.) To mention? (\textbf{\bn{উল্লেখ} \bn{করা}}, \emph{ullekh kôrā}.) Tension? (\textbf{\bn{উত্তেজনা}}, \emph{uttejônā}.) To spread a rumour? (\textbf{\bn{রটানো}}, \emph{rôṭāno}.) An uproar? (\textbf{\bn{হইচই}}, \emph{hôichôi}.) To whisper? (\textbf{\bn{ফিসফিস} \bn{করা}}, \emph{phisphis kôrā}.) A sensation? (\textbf{\bn{চাঞ্চল্য}}, \emph{chānchôlyô}.) A broadcast? (\textbf{\bn{সম্প্রচার}}, \emph{sômprôchār}.) A channel? (\textbf{\bn{চ্যানেল}}, \emph{chænel}.) A presenter? (\textbf{\bn{উপস্থাপক}}, \emph{upôsthāpôk}.) Live? (\textbf{\bn{সরাসরি}}, \emph{sôrāsôri}.) A viewer? (\textbf{\bn{দর্শক}}, \emph{dôrshôk}.)
\end{tcolorbox}
